% p221_1; source PDF page 221; manually verified editable table.
\begin{tikzpicture}[x=1cm,y=-1cm,inner sep=0pt,outer sep=0pt]
\path[use as bounding box] (0,0) rectangle (16.000,10.200);
\draw[line width=.4pt] (0.000,0.000) rectangle (4.500,0.900);
\node[align=center,text width=4.240cm,font={\linespread{1}\fontsize{10}{13.50}\selectfont\bfseries}] at (2.250,0.450) {行为类型};
\draw[line width=.4pt] (4.500,0.000) rectangle (9.700,0.900);
\node[align=center,text width=4.940cm,font={\linespread{1}\fontsize{10}{13.50}\selectfont\bfseries}] at (7.100,0.450) {结论};
\draw[line width=.4pt] (9.700,0.000) rectangle (16.000,0.900);
\node[align=center,text width=6.040cm,font={\linespread{1}\fontsize{10}{13.50}\selectfont\bfseries}] at (12.850,0.450) {范例};
\draw[line width=.4pt] (0.000,0.900) rectangle (4.500,5.200);
\node[align=left,text width=4.240cm,font={\linespread{1}\fontsize{10}{13.50}\selectfont}] at (2.250,3.050) {教唆或帮助他人自杀（没有欺骗或强迫他人自杀，没有支配力）};
\draw[line width=.4pt] (4.500,0.900) rectangle (9.700,5.200);
\node[align=left,text width=4.940cm,font={\linespread{1}\fontsize{10}{13.50}\selectfont}] at (7.100,3.050) {观点一：无罪。理由：坚持共犯从属性，实行者（自杀者）无罪，教唆者、帮助者也无罪。\\观点二（多数观点）：有罪。理由：生命要绝对保护，共犯从属性可以有例外};
\draw[line width=.4pt] (9.700,0.900) rectangle (16.000,5.200);
\node[align=left,text width=6.040cm,font={\linespread{1}\fontsize{10}{13.50}\selectfont}] at (12.850,3.050) {例1，甲教唆精神正常的成年人乙自杀，乙便自杀。\\例2，乙要自杀，甲提供毒药，乙自行喝下};
\draw[line width=.4pt] (0.000,5.200) rectangle (4.500,7.600);
\node[align=left,text width=4.240cm,font={\linespread{1}\fontsize{10}{13.50}\selectfont}] at (2.250,6.400) {教唆或帮助他人自伤（没有欺骗或强迫他人自伤，没有支配力）};
\draw[line width=.4pt] (4.500,5.200) rectangle (9.700,7.600);
\node[align=left,text width=4.940cm,font={\linespread{1}\fontsize{10}{13.50}\selectfont}] at (7.100,6.400) {无罪。理由：坚持共犯从属性};
\draw[line width=.4pt] (9.700,5.200) rectangle (16.000,7.600);
\node[align=left,text width=6.040cm,font={\linespread{1}\fontsize{10}{13.50}\selectfont}] at (12.850,6.400) {甲教唆乙自残，然后骗取保险金。乙照办。甲不构成故意伤害罪的教唆犯，但构成保险诈骗罪的教唆犯};
\draw[line width=.4pt] (0.000,7.600) rectangle (4.500,10.200);
\node[align=left,text width=4.240cm,font={\linespread{1}\fontsize{10}{13.50}\selectfont}] at (2.250,8.900) {乙承诺让甲杀乙，或重伤乙};
\draw[line width=.4pt] (4.500,7.600) rectangle (9.700,10.200);
\node[align=left,text width=4.940cm,font={\linespread{1}\fontsize{10}{13.50}\selectfont}] at (7.100,8.900) {承诺无效，甲构成故意杀人罪的实行犯，或者故意伤害罪的实行犯};
\draw[line width=.4pt] (9.700,7.600) rectangle (16.000,10.200);
\node[align=left,text width=6.040cm,font={\linespread{1}\fontsize{10}{13.50}\selectfont}] at (12.850,8.900) {［注意］被害人承诺的案件中，被害人不是实行者。在教唆、帮助他人自杀、自伤的案件中，被害人是实行者};
\end{tikzpicture}
